% A5 組版 (luatexja + ltjsbook) 用のプリアンブル。
% pandoc/build.sh が --include-in-header で読み込む。

\usepackage[most]{tcolorbox}
\usepackage{xcolor}

\definecolor{columnframe}{HTML}{4C5FD5}
\definecolor{columnback}{HTML}{F5F6FA}
\definecolor{noteframe}{HTML}{3B7DD8}
\definecolor{tipframe}{HTML}{8B5CF6}
\definecolor{cautionframe}{HTML}{D97706}
\definecolor{dangerframe}{HTML}{DC2626}

% 「実装から見ると」コラム。引数はバージョン (例: Unbound 1.26.0)。
% 読み飛ばせる塊として、本文から独立した枠にする (book.md §4.3)。
\newtcolorbox{implcolumn}[1]{%
  enhanced,
  breakable,
  colback=columnback,
  colframe=columnframe,
  coltitle=white,
  fonttitle=\bfseries\small,
  fontupper=\small,
  boxrule=0.4pt,
  leftrule=2pt,
  arc=1mm,
  top=4pt, bottom=4pt,
  title={実装から見ると\hfill\texttt{\footnotesize #1}},
}

% 注意書き。引数はタイトル。
\newtcolorbox{bookaside}[2]{%
  enhanced,
  breakable,
  colback=white,
  colframe=#1,
  coltitle=#1,
  colbacktitle=white,
  fonttitle=\bfseries\small,
  fontupper=\small,
  boxrule=0.4pt,
  leftrule=2pt,
  arc=0.5mm,
  top=4pt, bottom=4pt,
  title={#2},
}

\newenvironment{bookasidenote}[1]{\begin{bookaside}{noteframe}{#1}}{\end{bookaside}}
\newenvironment{bookasidetip}[1]{\begin{bookaside}{tipframe}{#1}}{\end{bookaside}}
\newenvironment{bookasidecaution}[1]{\begin{bookaside}{cautionframe}{#1}}{\end{bookaside}}
\newenvironment{bookasidedanger}[1]{\begin{bookaside}{dangerframe}{#1}}{\end{bookaside}}
